%%%%%%%%%%%%%%%%%%%%%%%%%%%%%%%%%%%%%%%%%%%%%%%%%%%%%%%%%%%%
%% CHAPTER 3 -- Networks of Simple Decisions
%%%%%%%%%%%%%%%%%%%%%%%%%%%%%%%%%%%%%%%%%%%%%%%%%%%%%%%%%%%%

\chapter{Networks of Simple Decisions}
\label{ntg:ch3}

\section*{Setting the Scene}
\addcontentsline{toc}{section}{Setting the Scene}

\subsection*{The problem}

We have, by now, a sense of what we want a Large Language Model
to do (Chapter~1) and a sense of how the model is improved over
time (Chapter~2). What we have not yet looked at is what is
\emph{inside} the model. The training procedure of the previous
chapter adjusts a great many internal numbers; we have not said
how those numbers are arranged, what they do when the model is
asked a question, or why a particular arrangement of them
should be capable of representing language at all.

This chapter is about the arrangement. It is about the artefact
the field calls a \emph{neural network} --- a name that is half
honest. The pieces of a neural network are loosely modelled on
the pieces of a biological brain, in the sense that they
combine many small inputs into a single output and pass the
output along to other pieces. They are not modelled on biological
brains in any deeper sense. Calling them ``neurons'' is a
historical accident that has outlasted its usefulness, and you
should think of the units as small mathematical decision-makers
rather than as anything resembling the cells in your skull.

\subsection*{How we got here}

The idea that a network of simple computing units could learn a
useful function from examples is older than the computers that
could run it. Warren McCulloch, a neurophysiologist, and Walter
Pitts, a young logician, published a paper in 1943 that
described how a network of idealised on-off units could, in
principle, compute any function expressible in symbolic logic.
The paper is now mostly read for its historical importance
rather than its mathematics. It established the question:
could a system of small interacting units, each doing something
trivial, collectively do something interesting?

Frank Rosenblatt, at Cornell in 1958, built the first working
answer. The \emph{perceptron} was a single computing unit that
took a list of inputs, multiplied each by a weight, added them
up, and produced an on-or-off output depending on whether the
total exceeded a threshold. The trick was that the weights
could be adjusted by showing the perceptron labelled examples:
if it got an example wrong, the weights were nudged in the
direction that would have got it right. Rosenblatt's machine
could, with enough training, learn to distinguish handwritten
shapes, simple visual patterns, and a small library of other
useful tasks. The press of the day was extravagant. The
\emph{New York Times} reported, with characteristic restraint,
that the perceptron was the embryo of a machine that would walk,
talk, see, write, and reproduce itself.

In 1969 Marvin Minsky and Seymour Papert published \emph{Perceptrons},
a book whose mathematical chapters proved that a single perceptron
could not represent certain very simple functions --- the most
famous being the exclusive-or, the operation that returns true
when exactly one of two inputs is true. Their proof did not say
that a network of \emph{many} perceptrons, arranged in layers,
could not do these things. But the field at the time did not
know how to train such networks, and the practical effect of the
book was to make funding for neural network research dry up.
The period that followed is now called the first AI winter, and
it lasted roughly until the mid-1980s.

The thaw came with backpropagation. Paul Werbos, in his 1974
Harvard PhD thesis, wrote down a recognisable derivation of the
algorithm by which a network of \emph{many} layers could be
trained --- the chain rule of calculus, applied carefully, will
tell you how each weight in each layer should be adjusted to
reduce the error at the output. Werbos's thesis sat largely
unread. The same algorithm was rediscovered, independently or
semi-independently, by several people in the mid-1980s, and it
was David Rumelhart, Geoffrey Hinton, and Ronald Williams who
published the version, in \emph{Nature} in 1986, that the field
finally noticed. The 1986 paper showed that a network of two or
three layers, trained by backpropagation, could solve problems
that the single-layer perceptron could not. Hinton would later
joke, with some justice, that the field had spent fifteen years
not knowing how to train two-layer networks because nobody
remembered Werbos.

The second renaissance ran from 1986 through the early 2000s.
Yann LeCun, at AT\&T Bell Labs, applied backpropagation to
networks for reading handwritten zip codes; the United States
Postal Service eventually deployed his system to sort about a
tenth of all American mail. Sepp Hochreiter and J\"urgen
Schmidhuber, in 1997, identified and largely solved the
\emph{vanishing gradient problem} that had been preventing
networks from learning across long sequences. By the late
1990s neural networks were a working tool in several specialist
domains. They were not yet the dominant tool in any of them.

The pivotal year was 2012. Alex Krizhevsky, Ilya Sutskever, and
Geoffrey Hinton entered a deep neural network into the ImageNet
visual recognition competition and won by a margin so large
that the competition organisers initially suspected an error.
The network was eight layers deep --- modest by today's
standards --- and it ran on consumer graphics cards rather than
on specialist hardware. The combination of a working algorithm
(backpropagation), enough training data (ImageNet's million
labelled images), and enough cheap computation (the GPU) had
converged. The field has not stopped accelerating since.

\subsection*{Where we stand}

A modern neural network used inside a Large Language Model is a
tall stack of small mathematical operations. Each layer of the
stack takes a list of numbers as input, multiplies the list by a
table of weights to produce a new list, adds a small adjustment
to each output, and then passes the result through a simple
nonlinear function (a function that, for example, replaces every
negative number with zero). The result is a new list of
numbers, of comparable size to the input, which becomes the
input to the next layer.

A typical Large Language Model today has between twenty and a
hundred such layers stacked on top of one another. The stack is
deeper than the deepest practical network of even ten years ago,
and the engineering tricks that made the depth trainable ---
careful initial values for the weights, normalisation of the
intermediate signals, and, crucially, the \emph{residual}
shortcut that lets each layer add its contribution to a running
total rather than transform the running total wholesale --- are
each named contributions of named researchers in the period from
2010 through about 2017.

The Transformer of Chapter~5 is, in essence, this stack of
layers with one further trick on top: each layer does not just
compute a function of its input; it computes a function of its
input \emph{and} of an attention-weighted summary of all the
other inputs in the sequence. We will see in Chapter~5 what
that means in detail. For now, the picture is enough: a deep
stack of small operations, each one slightly transforming a
list of numbers, with adjustments that the training procedure
slowly tuned.

\subsection*{What goes wrong}

The failure modes of neural networks are part of the
engineering vocabulary of the field, and they recur often
enough that even a non-technical reader benefits from
recognising them.

\emph{Vanishing} or \emph{exploding} signals: in a network of
the wrong design, the numbers passed from one layer to the next
either shrink to nothing or grow without bound, and the
training procedure either learns nothing at all or produces
gibberish. The fix is the careful initialisation, normalisation,
and residual-shortcut combination just described, and the
modern recipe for stacking layers is well understood enough
that this class of failure rarely happens in production.

\emph{Brittleness} to inputs unlike the training data: a network
trained on photographs from one camera will be measurably
worse on photographs from a different one. A model trained on
fluent English will be noticeably worse on text written by
someone for whom English is a second language. The training
process bakes the distribution of the training data into the
network in a way that cannot easily be removed.

\emph{Adversarial examples}: in the early 2010s, researchers
discovered that an image of a panda, modified by adding a
carefully chosen pattern of noise that was invisible to a
human viewer, could cause a state-of-the-art image classifier
to confidently label it a gibbon. The same phenomenon exists
for language models, in slightly different form: a prompt
modified by a few characters can produce qualitatively
different output. We will see this in Chapter~5 as
\emph{prompt brittleness} and in Chapter~12 as
\emph{prompt injection}.

\section{The idea, in plain language}
\addcontentsline{toc}{section}{The idea, in plain language}

\subsection*{What a single ``neuron'' actually does}

The smallest unit of a neural network is the artificial
``neuron'' that descends from Rosenblatt's perceptron. It does
something a committee chair would recognise. It takes a list of
inputs (think: the numerical opinions of various committee
members). It assigns each input a \emph{weight} (some members
have a stronger voice than others). It adds the weighted
opinions together. It adds a small \emph{bias} (the chair's
own predisposition). It then passes the total through a
\emph{nonlinear function} that decides what the unit's own
opinion will be --- often something as simple as ``if the total
is negative, output zero; if positive, output the total
unchanged''.

That is, in essence, the entire computation a single artificial
neuron performs. Several thousand of them in a row constitute a
\emph{layer}. Several layers stacked on one another constitute
a \emph{network}. The training of Chapter~2 adjusts the weights
and the biases. Nothing else in the network changes during
training.

The remarkable thing about this construction --- and it took the
field decades to come to peace with how remarkable it is --- is
that a network of these very simple decision-makers, given
enough of them and enough training, can represent extraordinarily
complicated functions. There are theorems that say so (the
universal approximation theorems of George Cybenko in 1989 and
Kurt Hornik in 1991, which proved that a sufficiently wide
single-layer network can approximate any reasonably well-behaved
function to any desired accuracy). The theorems do not, however,
say that the training procedure of Chapter~2 will reliably
\emph{find} such a representation, and most of the engineering
of modern neural networks is the work of arranging things so
that the training procedure does find one.

\subsection*{Why depth matters}

If the universal approximation theorems already say that a wide
enough single-layer network can represent any function, why does
the field bother with deep networks?

The answer is that ``can represent'' and ``can reasonably be
trained to represent'' are different things. A single-layer
network of the size required to represent a useful language
model would have so many neurons in its single layer that it
could not fit on any plausible amount of hardware, would
require an absurd amount of training data, and would in practice
fail to learn anything coherent. A deep network of more modest
width, on the other hand, can express the same function with
many fewer total parameters, because each layer can compose
the work of the layers below it.

The picture worth holding onto is a recipe for a recipe.
Imagine that the bottom layer of the network learns to detect
very simple features of the input --- in a vision network, the
presence of edges; in a language model, the presence of common
word fragments. The next layer learns to combine those simple
features into slightly more complex ones --- in vision, corners
and curves; in language, common multi-word phrases. The layer
above that combines the slightly-complex features into
moderately-complex ones --- shapes, syntactic structures. By
the time we are at the top of a fifty-layer network, the
features being combined are abstract enough to be useful for
the prediction task. No human has told the network what features
to detect at any layer. The training has worked them out.

This composition picture is partly right and partly a useful
fiction. We have decades of evidence from \emph{interpretability}
research that early layers in vision networks really do detect
edges, and that intermediate layers really do detect more
complex shapes. The picture is murkier for language models,
which have been studied less and which appear to organise their
internal computations along different lines than vision
networks do. But the basic intuition --- that depth lets the
network reuse the work of earlier layers when computing later
layers --- is the right one.

\subsection*{How the weights are adjusted: what backpropagation
actually does}

The picture of a deep stack of weighted-sum-and-nonlinearity
layers raises an obvious question: how does the training
procedure of Chapter~2 know which of the millions of weights
to nudge, and in which direction?

The answer is \emph{backpropagation}. The name is a contraction
of ``backward propagation of errors'', which is a reasonably
accurate description of what it does. Training a network on a
single example proceeds in two passes.

The first pass --- the \emph{forward pass} --- sends the input
through the network, layer by layer, until the network produces
a prediction. The prediction is compared against the correct
answer. A number is computed that describes how wrong the
prediction was. This number is the \emph{loss}.

The second pass --- the \emph{backward pass} --- sends the
error back through the network in reverse. For each weight,
the backward pass computes how much the loss would change if
that weight were nudged very slightly upward. This is a
direction: nudge the weight in the opposite direction, and
the loss should fall. The key insight is that these individual
contributions can be computed efficiently, even for a network
of a hundred layers, because the loss is a composition of
functions, and there is a rule in calculus --- the chain rule
--- that tells you exactly how to decompose the effect of any
weight in the chain onto the final output. Backpropagation
is, at bottom, the chain rule applied carefully and in reverse.

This is why training is expensive. For every example in the
training set, you perform a forward pass and then a backward
pass. A backward pass costs roughly as much as a forward pass.
And this must be done for every weight, across every example,
for many thousands of iterations before the loss converges.
The cost is why training a frontier model requires weeks on
tens of thousands of specialised processors.

This is also where GPUs enter the picture. Within a single
layer, the forward-pass multiplications --- each neuron
computing its weighted sum independently of the others in the
same layer --- are independent. They do not wait for each
other. A graphics card carries thousands of small processors
that can perform these independent multiplications in parallel,
which is why a task originally designed for rendering video
games turned out to accelerate neural-network training by two
or three orders of magnitude.

Finally, this is where the vanishing gradient problem lives.
The chain rule, applied across many layers, involves
multiplying many small numbers together. A product of, say,
fifty numbers each slightly less than one is a number very
close to zero. When the gradient reaching the early layers is
very close to zero, those weights receive almost no update
signal, and the network's early layers fail to learn. The
residual shortcut mentioned earlier --- the trick of adding
each layer's output to the running total rather than replacing
it --- short-circuits this multiplication and lets the gradient
flow through the addition rather than through the chain.
It is the fix that made networks of fifty or a hundred layers
trainable at all.

\subsection*{What the network is not}

Several things follow from the picture we have just laid out
that are worth stating explicitly because they tend to be
muddled in popular discussion.

The network does not contain a database. It does not look
things up. It is a function from inputs to outputs, computed by
performing arithmetic on the inputs. The arithmetic is
parameterised by the weights, and the weights were set by the
training procedure of Chapter~2. There is no place inside the
network where ``what the model knows'' is written down in
retrievable form. There is only the procedure that, when the
prompt is the input, produces a particular distribution over
next words as the output. The distinction between this and
having a database matters when we discuss retrieval in
Chapter~10 --- retrieval, in essence, gives the model the
database it does not have.

The network does not reason in the sense a human does. It
performs arithmetic. The arithmetic was tuned, by training,
to produce outputs that resemble the outputs a corpus of
human-written text produced. When the model generates a
chain of reasoning that arrives at a correct answer, what has
happened is that the corpus contained many similar chains of
reasoning, and the trained model has learned to imitate them.
This is not the same as having reasoned. It is, often,
indistinguishable from having reasoned at the level of the
output, which is part of what makes the technology so
disconcerting.

The network is not, in any technical sense, a black box ---
though it is fashionable to call it one. Every operation
inside the network is exposed to inspection. Researchers can,
and do, study what each neuron in each layer responds to.
What is true is that the number of operations is so large
that no human can summarise them all, and that the
interpretations researchers extract from individual neurons
are post-hoc reconstructions rather than authoritative
descriptions of what the network is ``doing''. The honest
phrase is that the network is \emph{difficult to interpret},
not that it is opaque in principle.

\section{The engineering consequence}
\addcontentsline{toc}{section}{The engineering consequence}

\paragraph{The function class is enormously powerful but
uninterpretable.} A network of the size used in modern Large
Language Models can represent essentially any function for which
sufficient training data exists. This is its strength: the
research community does not have to design an architecture
specific to each task. It is also its weakness: the network's
representation of a particular function is a tangle of weights
that no human can read directly. The practical consequence is
that \emph{interpretability} --- the discipline of working out
what a network is doing internally --- is an active research
area whose results are partial, post-hoc, and contested.

\paragraph{The same architecture can be trained for many
different tasks.} The deep stack of layers we have described
is, with minor variations, the same architecture used to
identify objects in photographs, transcribe speech, predict
protein folding, and generate language. The differences are in
the input format, the output format, and the training data.
This is why a research advance in one application area --- a
better way of stacking layers, say, or a more efficient
optimiser --- tends to propagate quickly to all the others. It
is also why the failure modes we will discuss for Large
Language Models have analogues in image recognition, speech,
and beyond.

\paragraph{Training and serving have very different costs.}
Training a neural network involves running each example forward
through the network, computing the error, and then running the
error backward through the network to update the weights. The
backward pass costs roughly as much as the forward pass.
Serving the trained network --- the operation that happens
when a user sends a query and waits for the response --- only
requires the forward pass. The serving cost per query is
therefore typically a few thousandths to a few cents of a
trained model, while the training cost is in the tens or
hundreds of millions of dollars. The lopsided economics drive
much of how the industry is organised.

\paragraph{Deeper networks are not always better networks.}
There is a long-standing temptation, in this field, to assume
that adding layers will improve a model's performance. The
truth is more subtle. Adding layers helps until the additional
layers run out of new things to compute, or until the
additional weights run out of training signal to learn from.
The right depth for a model depends on the task, on the data,
on the budget, and on a number of other choices that
interact. We will see in Chapter~7 that the relationship
between depth, width, training data, and capability is one of
the things the scaling laws were created to make precise.

\section{What goes wrong, and why}
\addcontentsline{toc}{section}{What goes wrong, and why}

The failure modes of neural networks are largely failures of
mismatch: the network was trained on one distribution, the
network is asked to produce an answer about a different one, and
the network's behaviour at the boundary is unpredictable.

\paragraph{Brittleness near the edge of training distribution.}
The network's behaviour on inputs very close to its training
data is, on the whole, reliable. Its behaviour on inputs that
are unlike anything in its training data is unpredictable in
ways that the network itself does not flag. There is, in
particular, no built-in mechanism by which the network reports
``I have not seen anything like this question before'' --- the
network has only the same response mechanism for known and
unknown inputs, and it produces a confident-looking output in
both cases. Building systems that detect ``out-of-distribution''
inputs is an open problem, and the partial solutions we have
are usually external to the model rather than inside it.

\paragraph{The black-box framing, reconsidered.} It is common
to describe neural networks as black boxes whose decisions
cannot be explained. The framing is partly accurate and partly
misleading. It is accurate in the sense that no human can
follow, in detail, every operation a billion-parameter model
performs on a particular input. It is misleading in the sense
that researchers can and do reverse-engineer significant parts
of what models compute --- a discipline called \emph{mechanistic
interpretability} that has, in the last few years, produced
recognisable accounts of how certain classes of computation are
performed in certain models. A more honest framing is that
neural networks are difficult to interpret rather than
impossible to interpret, and that the difficulty is a real
constraint on regulation, audit, and trust.

\paragraph{Adversarial inputs.} A small, deliberately chosen
perturbation of an input can cause the model to produce a
qualitatively different output. In vision models this takes the
form of imperceptible pixel patterns that change a
classification. In language models it takes the form of
sequences of characters that look like noise to a human reader
but that override the model's normal behaviour --- the technical
name is a \emph{jailbreak}. The vulnerability is intrinsic to
the function-approximation framing of neural networks: a
network that interpolates smoothly between its training
examples will also interpolate into regions of input space where
its training was sparse, and a clever attacker can find those
regions.

\paragraph{The training-distribution prison.} A model trained
on a particular corpus will, no matter how well trained, behave
with the assumptions and biases of that corpus. A model trained
predominantly on English text will treat English as the default;
a model trained predominantly on the public web will inherit
the public web's distribution of opinions, errors, and
omissions. There is no architectural fix for this. The fix is
either better training data (which is expensive and contested),
retrieval (Chapter~10), or careful prompt engineering at the
edge (which is brittle).

\section*{Bridges}
\addcontentsline{toc}{section}{Bridges}

This chapter built on the surprise-minimisation picture of
Chapters~1 and~2 and gave us the function class --- the deep
stack of weighted-sum-and-nonlinearity layers --- that the
training procedure of Chapter~2 is adjusting. Every chapter
that follows assumes this picture. Chapter~4 puts the
\emph{embeddings} that turn words into lists of numbers at the
bottom of the stack. Chapter~5 adds the \emph{attention}
mechanism, the trick that turned ordinary deep networks into
the Transformer. Chapter~7 returns to the question of how
deep, how wide, and on how much data, framed as a scaling
problem. Chapter~8 returns to the question of how to modify a
trained network without retraining it from scratch. The next
chapter --- \emph{From Letters to Vectors} --- explains how the
text on your screen ever becomes a list of numbers in the first
place.
